\begin{tikzpicture}
  % zip łączy elementy o tych samych indeksach w pary
  \def\sx{1.25}
  \foreach \v [count=\i from 0] in {5,3,7,2} {
    \node[komorka] (a\i) at (\i*\sx, 0) {\v};
    \node[indeks, above=0.4mm of a\i] {\i};
  }
  \foreach \v [count=\i from 0] in {1,-2,3} {
    \node[komorka] (b\i) at (\i*\sx, -2.2) {\v};
    \node[indeks, below=0.4mm of b\i] {\i};
  }
  \node[komorka szara] at (a3) {2};
  \node[kod, anchor=east] at (-0.7, 0) {a};
  \node[kod, anchor=east] at (-0.7, -2.2) {b};
  \foreach \i/\p in {0/{(5, 1)}, 1/{(3, -2)}, 2/{(7, 3)}} {
    \draw[draw=akcent, thick] (a\i.south) -- (b\i.north);
    \node[font=\ttfamily\scriptsize, fill=white, inner sep=1.5pt, text=akcent] at ($(a\i.south)!0.5!(b\i.north)$) {\p};
  }
  \node[notka, anchor=west, text width=34mm] at ($(a3.east)+(0.25,-1.1)$) {\kod{2} nie ma pary —\\\kod{zip} kończy na\\\textbf{krótszej} liście};
  \draw[->, draw=szary, densely dashed] ($(a3.east)+(0.3,-0.75)$) to[bend right=20] ($(a3.south)+(0.1,-0.05)$);
  % kod
  \node[notka, anchor=west, text width=60mm] (code) at (7.6, -1.1) {%
    \kod{\textcolor{fiolet}{for} x, y \textcolor{fiolet}{in} zip(a, b):}\\
    \kod{\ \ \ \ print(x, y)}\\
    \textcolor{szary}{\kod{\# 5 1 / 3 -2 / 7 3}}\\[2mm]
    \kod{\textcolor{fiolet}{for} i, x \textcolor{fiolet}{in} enumerate(a):}\\
    \kod{\ \ \ \ print(i, x)}\\
    \textcolor{szary}{\kod{\# 0 5 / 1 3 / 2 7 / 3 2}}};
  \begin{scope}[on background layer]
    \node[ramka, fit=(code), inner sep=4pt] {};
  \end{scope}
\end{tikzpicture}
